% TRANSLATION-PRODUCER DRAFT — UNCHECKED
% Work: Emmy Noether, Paper 41, Korean unit U12
% Source custody: ${PUBLIC_INTERLANGUAGE_ROOT}\03_projects\language_management\cjk\03_working_translations\noether_paper41_zh_translation_001_20260722\source\Noether_Paper41_CurrentGermanAuthority_interval.tex
% Source locator: interval lines 145--151
% Source unit bytes (LF-normalized UTF-8): 1327
% Source unit SHA-256: D0259D62D1D60476EE3F16E2E22BDFB8CDAD7A4978DC12683A5C0999AEA4DDBE
% Source snapshot SHA-256: C265058425E5E2D1A2289CC03A9DDEDDDF4803A3215DC3F173B93E7AB69D60ED
% Historical whole-source SHA-256: 443EF950D7D45DC6D9E44A9B87501620C10DA873E50E5F2B253ECCAE6A946D27
% Producer boundary: Korean translation only. No source/scan comparison, Korean or formula review,
% compilation, rendering, assembly, packaging, certification, approval, or self-check was performed.

두 보조정리로부터 \emph{주종정리의 증명}이 곧바로 따라 나온다. 각 \(\mathfrak c_S\)가 류 \(\bar c_S\)에 속하는 하나의 아이디얼을 뜻한다고 하자. 그러면 류들에 관한 가정은---인자계의 유도된 류분할의 정의를 고려하면---\(\mathfrak c_S\)의 변환량이 보조정리 2의 가정을 만족하는 주아이디얼 \((a_{S,T})\)이 된다는 바로 그 사실을 말한다. 따라서 주아이디얼 \((d_S)\)를 적절히 택하여 \(\mathfrak d_S=\mathfrak c_S/(d_S)\)로 둔 아이디얼들이 여전히 류 \(\bar c_S\)에 속하면서 보조정리 1의 가정을 만족하게 할 수 있다. 즉 \(\mathfrak d_S\)의 변환량은 단위아이디얼과 같아진다. 그러므로 보조정리 1에 따라 \(\mathfrak d_S=\mathfrak b^{1-S}\)를 만족하는 아이디얼 \(\mathfrak b\)가 존재한다. 이는 류 \(\bar c_S\)가 류 \(\bar b\)의 기호적 \((1-S)\)제곱이 된다는 뜻이다.

순환 경우의 주종정리가 이미 알려진 정리로 귀착된다는 사실은 다음으로부터 따른다. 즉 정규화를 기초로 삼으면, 유도된 류분할의 주류는 모든 분기소에서 \(\alpha\)가 노름잉여가 되는 주아이디얼 \((\alpha)\)들의 군으로 귀착된다.

\begin{center}
(1932년 10월 27일 접수.)
\end{center}
